%% LyX 2.4.4 created this file.  For more info, see https://www.lyx.org/.
%% Do not edit unless you really know what you are doing.
\documentclass[oneside,english]{book}
\usepackage[T1]{fontenc}
\usepackage[latin9]{inputenc}
\setcounter{secnumdepth}{3}
\setcounter{tocdepth}{3}
\usepackage{units}
\usepackage{amsmath}
\usepackage{amsthm}
\usepackage{amssymb}
\usepackage{cancel}
\usepackage{geometry}
\geometry{verbose,tmargin=1cm,bmargin=1cm,lmargin=1.5cm,rmargin=1.5cm}
\usepackage{wasysym}

\makeatletter
%%%%%%%%%%%%%%%%%%%%%%%%%%%%%% Textclass specific LaTeX commands.
\theoremstyle{definition}
\newtheorem{defn}{\protect\definitionname}
\theoremstyle{definition}
\newtheorem{example}{\protect\examplename}
\theoremstyle{definition}
\newtheorem{xca}{\protect\exercisename}
\theoremstyle{plain}
\newtheorem{thm}{\protect\theoremname}
\ifx\proof\undefined\
  \newenvironment{proof}[1][\proofname]{\par
    \normalfont\topsep6\p@\@plus6\p@\relax
    \trivlist
    \itemindent\parindent
    \item[\hskip\labelsep
          \scshape
      #1]\ignorespaces
  }{%
    \endtrivlist\@endpefalse
  }
  \providecommand{\proofname}{Proof}
\fi
\theoremstyle{plain}
\newtheorem{lem}{\protect\lemmaname}
\theoremstyle{plain}
\newtheorem{prop}{\protect\propositionname}

%%%%%%%%%%%%%%%%%%%%%%%%%%%%%% User specified LaTeX commands.
\usepackage{hyperref}
\usepackage[usenames,dvipsnames]{xcolor} %adds additional color options
\hypersetup{colorlinks=true, urlcolor=Blue}%Blue

\usepackage{multicol}

\usepackage{tikz}
\usepackage{tikz-3dplot}
\usetikzlibrary{calc,arrows,shapes,backgrounds,matrix,shapes.geometric,fit,trees,positioning}

\usetikzlibrary{patterns}
\usetikzlibrary{plotmarks}
\usepackage{pgfplots}
\pgfplotsset{compat=newest}

\usepackage{calc}
\usepackage[nomessages]{fp}

\usepackage{datetime}

%%%%%commands
\newcommand{\mb}[1]{$\mathbb{#1}$}
\newcommand{\funct}[2]{$f\colon#1\rightarrow#2$}
% Added by lyx2lyx
\setlength{\parskip}{\smallskipamount}
\setlength{\parindent}{0pt}

\makeatother

\usepackage{babel}
\providecommand{\definitionname}{Definition}
\providecommand{\examplename}{Example}
\providecommand{\exercisename}{Exercise}
\providecommand{\lemmaname}{Lemma}
\providecommand{\propositionname}{Proposition}
\providecommand{\theoremname}{Theorem}

\begin{document}

\chapter{Integers and Congruence\label{chap:Congruence}}

\section{Integers}

Much of abstract algebra studies the set of integers, subsets of integers,
or sets which are equivalent to subsets of integers. Thus to move
forward, it will behoove us to step back and investigate the properties
of integers. You have been using many of these properties all of your
life. Taking a rigorous look at these properties will not only enrich
your understanding, but will allow you to apply these fundamental
properties to more general sets. 
\begin{defn}
\textbf{The Well Ordering Principle. }Every nonempty set of positive
integers contains a smallest member. 
\end{defn}
We will treat the well ordering principle as an axiom.\footnote{Meaning we will just assume it to be true for the integers. The well
ordering principle sometimes given as a theorem proved with an induction
(see below). As it turns out, induction and the well ordering principle
are logically equivalent, and we will instead show that induction
follows from the well ordering principle.}
\begin{defn}
We say that a nonzero integer $d$ \emph{divides }an integer $m$
(or say that $d$ is a \emph{divisor of $m$}) if $m$ can be written
as, $m=dn$ for some integer $n$, denoted $d|m$.

\begin{itemize}
\item $d|m\Rightarrow m=dn$ for some integer $n$.
\item $d\cancel{|}m\Rightarrow m\neq dn$ for any integer $n$.
\end{itemize}
\end{defn}
\begin{example}
$3|12$ since $12=3\cdot4$ (which also shows that $4|12$). $3\cancel{|}10$
since there is no integer $n$ such that $3\cdot n=10$ (solving for
$n$, $n=\nicefrac{10}{3}\notin\mathbb{Z}$).
\end{example}
\begin{xca}
Prove that If $a|b$ and $b|c$ then $a|c$.
\end{xca}

\begin{xca}
If $a|b$ and $a|c$ then $a|b\pm c$.
\end{xca}
\begin{defn}
A prime $p$ is an integer greater that has no divisors other than
$1$ and $p$.

\begin{itemize}
\item Equivalently, $p$ is prime if and only if $d|p\Rightarrow d=1\mbox{ or }d=p$.
\end{itemize}
\end{defn}
\begin{example}
$5$ is prime since $1<5$ and $2,3,4\cancel{|}5$.
\end{example}
\begin{xca}
$\bigstar$ Prove that there are infinitely many prime numbers. 
\end{xca}
\begin{thm}
\textbf{\textup{Division Algorithm. }}\textup{\emph{Let $a,b\in\mathbb{Z}$
and $a,b>0$ then there exists a unique $q,r\in\mathbb{Z}$ such that
$a=bq+r$ where $0\leq r<b$.}}\emph{}

\begin{itemize}
\item \emph{$r$ is called the remainder upon dividing $a$ by $b$}. 
\end{itemize}
\end{thm}
\begin{example}
If you were to divide $13$ by $3$, the division algorithm tells
us that $13=3\cdot4+1$; where $a=13$, $b=3$, $q=4$, and $r=1$.
Once we have chosen $a$ and $b$ (or $q$), the theorem tells us
that $a$ can be written uniquely in this form. 
\end{example}
\begin{itemize}
\item Note: we can write $13$ as $13=4\cdot2+5$, $13=4\cdot1+9$, or $13=4\cdot5-7$.
However the remainders are $5,9>4$ in the former cases and, $-7<0$
in the latter case. The assumption that $0\leq r<b$ is what guarantees
uniqueness.
\end{itemize}

\begin{example}
$57=12\cdot4+9$. 

\begin{proof}
There are two parts to prove, existence and uniqueness. First we show
that you can write any integer $a$ as $bq+r$ where $0\leq r<b$.
Let $S=\{a-bq\colon q\in\mathbb{Z}\mbox{ and }a-bq\geq0\}$. First
note that $S$ is nonempty. If $a-b<0$ then let $k=-1$ so that $a-(-1)b>0$.
If $a-b\geq0$ then let $k=1$. 

If $0\in S$ then $a=bk$ which implies that $b|a$ and then we can
let $q=\nicefrac{a}{b}$ (an integer since $b|a$) and $r=0$ so that
$a=b(\nicefrac{a}{b})+0$. 

If $0\notin S$ then we can apply the well ordering principle (since
$S$ is nonempty) to choose the smallest number in $S$, say $r=a-bq$.
Then $a=bq+r$; $r\geq0$ per our definition of $S$ so we need only
show that $r<b$. If $r\geq b$ then $r-b\geq0\Rightarrow a-bq-b=a-b(q+1)$
so that $a-b(a+1)\in S$. But then $a-bq>a-b(q+1)$ and we already
assumed that $a-bq=r$ was the \emph{smallest }number in $S$ --a
contradiction. So $r<b$.

Now we show uniqueness. Let $a=bq+r$ and $a=bq'+r'$ such that $0\leq r<b$
and $0\leq r'<b'$ for integers $q,q',r,r'$. Without loss of generality\footnote{often abbreviated to WLOG, WOLOG or w.l.o.g.; less commonly stated
as without any loss of generality or with no loss of generality. WLOG
means that although a premise of the proof is being narrowed to a
special case, the proof of that special case can be easily applied
to every other case.} we can assume that $r'\geq r$. $bq+r=bq'+r'$ and $bq'+r'-(bq+r)=0\Rightarrow b(q'-q)=r'-r$.
Thus $b|(r'-r)$ and $0\leq r'-r\leq r'<b$. The only way $b$ can
divide a number smaller than itself (and still be an integer) is if
that number is $0$. So then $r'-r=0\Rightarrow r=r'$ which also
shows that $q'=q$.
\end{proof}
\end{example}
\begin{defn}
The \emph{greatest common divisor }(or GCD) of integers $a$ and $b$
is the largest integer that divides both $a$ and $b$; denoted $\gcd(a,b)$.
When $\gcd(a,b)=1$ we say that $a$ and $b$ are \emph{relatively
prime. }
\end{defn}
\begin{example}
$\gcd(10,6)=2$ since $10=2\cdot5$ and $6=2\cdot3$.
\end{example}

\begin{example}
$1100=2^{2}5^{2}11$ and $440=2^{3}5\cdot11$. So $\gcd(1100,440)=2^{2}5\cdot11=220$.
\end{example}
\begin{defn}
The \emph{lowest common multiple }(or LCM) of integers $a$ and $b$
is the smallest integer that is a multiple of both $a$ and $b$.
\end{defn}
\begin{example}
$10=5\cdot2$ and $6=2\cdot3$ so $\mbox{lcm}(10,6)=5\cdot2\cdot3=60$.
Similarly the $\mbox{lcm}(1100,440)=2^{3}5^{2}11=2200$.
\end{example}
\begin{xca}
Determine the following:

\begin{enumerate}
\item $\gcd(2^{4}3^{2}5\cdot7^{2},2\cdot3^{3}7\cdot11)$
\item The lowest common multiple: $\mbox{lcm}(2^{4}3^{2}5\cdot7^{2},2\cdot3^{3}7\cdot11)$.
\end{enumerate}
\end{xca}

\begin{xca}
$\bigstar$ Suppose that $a\neq0$, $b\neq0$, and $\mbox{lcm}(a,b)=m$.
Prove that if $n$ is a common multiple of $a$ and $b$ then $m|n$. 
\end{xca}

\begin{xca}
$\bigstar$ Let $a\neq0$, $b\neq0$, and $d=\gcd\left(a,b\right)$.
Prove that if $d'$ is any common divisor of $a$ and $b$ then $d'|d$.
\end{xca}
The next theorem plays an important role in our study of algebra.
The proof relies on both the well ordering principle and the division
algorithm.
\begin{thm}
\textbf{\emph{\label{thm:GCD-is-a-linear-combo}GCD is a linear combination.
}}For any nonzero integers $a$ and $b$, there exists integers $s$
and $t$ such that $\gcd(a,b)=as+bt$. Furthermore, $\gcd(a,b)$ is
the smallest positive integer that can be written in the form $as+bt$.
\end{thm}
\begin{itemize}
\item In particular, $a$ and $b$ are relatively prime if and only if there
exists integers $s$ and $t$ such that $1=as+bt$.

\begin{proof}
Let $a,b$ be two nonzero integers. Let $S=\{am+bn\colon m,n\in\mathbb{Z}\mbox{ and }am+bn>0\}$.
Note that $S\neq\emptyset$ since if $am+bn\leq0$ then $a(-m)+b(-n)>0$.
By the well ordering principle we can select the smallest number in
$S$, say $d=as+bt$. We want to show that $d=\gcd(a,b)$. 

Using the division algorithm, write $a=dq+r$ where $0\leq r<d$.
If $r>0$ then $r=a-dq=a-(as+bt)q=a-asq-btq=a(1-sq)+b(-tq)\in S$.
But this contradicts the assumption that $d$ is smallest number in
$S$ since $a(1-sq)+b(-tq)<as+bt$. So then $r=0$ implying that $d|a$.
Similarly we can show that $d|b$. This shows that $d$ is a common
divisor of $a$ and $b$. 

Now we show that $d$ is the largest common divisor. Suppose that
$d'$ is another common divisor of $a$ and $b$. So then $a=d'x$
and $b=d'y$ for integers $x$ and $y$. Now $d=as+bt=(d'x)s+('dy)t=d'(xs+yt)$
which implies that $d'|d$. Therefore $d'<d$.
\end{proof}
\end{itemize}
\begin{example}
$\gcd(10,6)=2$ and $2=10(2)+6(-3)$. 
\end{example}

\begin{example}
$\gcd(4,15)=1$. So $4$ and $15$ are relatively prime. $1=4(4)+15(-1)$. 
\end{example}
We will see that theorem \ref{thm:GCD-is-a-linear-combo} is a very
useful theoretical tool. However it does not tell us how to actually
find the linear combinations. Fortunately, we have the Euclidean Algorithm\footnote{From proposition 2 of book seven of the \emph{Elements }written approximately
c.\,300 BC.}.
\begin{example}
\textbf{Euclidean Algorithm. }By repeatedly using the division algorithm
the $\gcd(a,b)$ can always be found by producing a decreasing sequence
of remainders. The algorithm is as follows:
\[
\begin{array}{rclcl}
a & = & bq_{1}+r_{1} & \mbox{ \,\,\,} & 0<r_{1}<b\\
b & = & r_{1}q_{2}+r_{2} &  & 0<r_{2}<r_{1}\\
r_{1} & = & r_{2}q_{3}+r_{3} &  & 0<r_{3}<r_{2}\\
 & \vdots\\
r_{n-3} & = & r_{n-1}q_{n-1}+r_{n-1} &  & 0<r_{n-1}<r_{n-2}\\
r_{n-2} & = & r_{n-1}q_{n}+r_{n} &  & 0<r_{n}<r_{n-1}\\
r_{n-1} & = & r_{n}q_{n+1}+0
\end{array}
\]
By the well ordering principle, this algorithm will eventually result
in a remainder of $0$. Then the last nonzero remainder, $r_{n}$,
is the greatest common divisor.

For example, to find the $\gcd(2520,154)$: 
\[
\begin{array}{rcl}
2520 & = & 154(16)+56\\
154 & = & 56(2)+42\\
56 & = & 42(1)+14\\
42 & = & 14(3)+0
\end{array}
\]
So $\gcd(2520,154)=14$. We can also use this algorithm to find a
linear combination of $2520$ and $154$ that sums to $14$. Using
the above results, we rewrite the remainders as follows: 
\[
\begin{array}{rclcrcl}
56 & = & 42(1)+14 & \mbox{\ensuremath{\Rightarrow}} & 14 & = & 56-42(1)\\
154 & = & 56(2)+42 & \Rightarrow & 14 & = & 56-[154-56(2)]\\
 &  &  &  & 14 & = & (3)56-154\\
2520 & = & 154(16)+56 & \Rightarrow & 14 & = & (3)[2520-154(16)]-154\\
 &  &  &  & 14 & = & 2520(3)+154(-49)
\end{array}
\]
\end{example}
\begin{xca}
Find integers $s$ and $t$ so that $1=7\cdot s+11\cdot t$. Show
that $s$ and $t$ are not unique. 
\end{xca}

\begin{xca}
Let $d=\gcd(a,b)$. If $a=da'$ and $b=db'$, show that $\gcd(a',b')=1$.
\end{xca}

\begin{xca}
$\bigstar$ If $c^{2}=ab$ and $\gcd\left(a,b\right)=1$, prove that
$a$ and $b$ are perfect squares (Hint: By the theorem above, $a$
and $b$ can be written as the product of primes. They are perfect
squares if and only if each prime in this prime factorization has
an even power). 
\end{xca}
\begin{lem}
\textbf{\emph{Euclid's Lemma. }}Let $p$ be a prime and $a$ and $b$
be integers. If $p|ab$ then $p|a$ or $p|b$.\footnote{This lemma first appeared as proposition $30$ in book seven of Euclid's
\emph{Elements}. }

\begin{proof}
Suppose that $p|ab$ but $p\cancel{|}a$. We want to show that $p|b$.
Since $p\cancel{|}a$, $\gcd(p,a)=1$ and by theorem \ref{thm:GCD-is-a-linear-combo}
we can write $1=sp+ta$. Multiplying both sides by $b$ we have $b=bsp+bta$.
Since $p|a$, $a=pn$ for some integer $n$, and 
\begin{eqnarray*}
b & = & bsp+bta\\
\Rightarrow b & = & bsp+btpn\\
\Rightarrow b & = & p(bs+btn)
\end{eqnarray*}
so $p|b$ as needed. 
\end{proof}
\end{lem}
\begin{itemize}
\item Euclid's lemma is not necessarily true when applied to non-primes.
For example, $10|25\cdot2$ but $10\cancel{|}25$ and $10\cancel{|}2$.
\end{itemize}

\section{\textit{\emph{Modular Arithmetic}}}

Modular arithmetic is a system of arithmetic where numbers ``wrap
around'' or ``reset'' after reaching a set value called the \emph{modulus}.
You have probably been using modulus arithmetic all your life without
giving it much thought. One common use is a 12-hour clock. 

If the time is 5:00 then 11 hours later we do not say it is 16:00
(unless you are in the military). Time ``wraps'' around every 12
hours -modulus 12. Using the division algorithm, $16=12(1)+4$, and
we have a remainder of $4$. $16$ ''wraps'' $4$ hours past $12$.
We write this as 
\[
16\equiv4\bmod12
\]
read ``$16$ is congruent/equivalent to $4$ mod $12$''. Adding
$12$ hours to $16$ will wrap around to the same number, $4$, as
will adding $24$, $36$, $48$ etc. 
\begin{eqnarray*}
4 & \equiv & 4\bmod12\\
16 & \equiv & 4\bmod12\\
28 & \equiv & 4\bmod12\\
40 & \equiv & 4\bmod12\\
\vdots\\
4+12(n) & \equiv & 4\bmod12
\end{eqnarray*}
So $4$, $16$, $28$, etc.\,are all members of an infinite set we
call a \emph{class}, denoted $[4]$.
\begin{defn}
Let $a$ and $n$ be integers with $n>0$. The \emph{congruence class
of a modulo n }(denoted $[a]$) is the set of integers that are congruent
to $a$ modulo $n$, i.e., 
\[
[a]=\left\{ b|b\in\mathbb{Z}\mbox{ and }b=a\bmod n\right\} 
\]
\end{defn}
\begin{example}
The congruence class of $3$ modulo $5$ is $[3]=\left\{ \dots,-2,3,8,13,\dots\right\} $.
\end{example}
There are several useful ways to define $a\equiv b\bmod n$. However
when first introduced to the subject it is often easiest to think
of modular arithmetic in terms of the remainder. 
\begin{itemize}
\item If $a=qn+r$ where $q$ is the \emph{quotient }and $r$ is the \emph{remainder
}upon dividing $a$ by the \emph{modulus }$n$, then $a\equiv r\bmod n$. 
\end{itemize}
\begin{example}
If $20$ is divided by $3$ then the remainder is $2$. So $20\equiv2\bmod3$
since $20=(6)3+2$. 
\end{example}

Of course if $20\equiv2\bmod3$ then one would think that $2\equiv20\bmod3$
(as congruence/equivalence should be an equivalence relation), however
dividing $3$ by $20$ is not permissible as any remainder must be
greater than $0$ but less than $20$. For this reason we define $a\equiv b\bmod n$
as follows:
\begin{defn}
$a\equiv b\bmod n$ if and only if $n|(a-b)$.
\end{defn}
\begin{example}
Using the definition above, $20\equiv2\bmod3$ since $3|(20-2)$ and
$3\equiv20\bmod3$ since $3|(2-20)$.
\end{example}

\begin{example}
If If $a=qn+r$ then $r=a-qn\Rightarrow qn=a-r$. So then $n|a-r$
and $a\equiv r\bmod n$.
\end{example}

\begin{example}
Examples of modular arithmetic:
\begin{eqnarray*}
7 & \equiv & 7\bmod13\mbox{\,\,\,\ since }7=(0)13+7\\
20 & \equiv & 7\bmod13\mbox{\,\,\,\ since \ensuremath{20=(1)13+7}}\\
7 & \equiv & 20\bmod13\,\,\,\mbox{since }7=(-1)13+7\\
48 & \equiv & 0\bmod12\,\,\,\mbox{ since }48=(4)12+0\\
75 & \equiv & 25\bmod50\,\,\,\mbox{ since }75=(1)50+25
\end{eqnarray*}
\end{example}
\begin{prop}
If $a\equiv b\bmod n$ and $c\equiv d\bmod n$ for integers $a,b,c,d,n$\footnote{This assumption is only necessary for part 2. While part 1 is true
for all real numbers part may fail for non-integers.} then 

\begin{enumerate}
\item $a+c\equiv(b+d)\bmod n$
\item $ac\equiv bd\bmod n$
\end{enumerate}
\begin{proof}
$a\equiv b\bmod n\Rightarrow n|(a-b)\Rightarrow a-b=k_{1}n$ and $c\equiv d\bmod n\Rightarrow n|(c-d)\Rightarrow k_{2}n$
for some integers $k_{1}$ and $k_{2}$. 

Part 1. $(a+c)-(b+d)=a+c-b-d=(a-b)+(c-d)=k_{1}n+k_{2}n=n(k_{1}+k_{2})$.
Thus $n|[(a+c)-(b+d)]$ which implies that $a+c\equiv(b+d)\bmod n$.

Part 2. $a-b=k_{1}n$ and $c-d=k_{2}n$ implies that $a=k_{1}n+b$
and $c=k_{2}n+d$. 
\begin{eqnarray*}
ac & = & (k_{1}n+b)(k_{2}n+d)\\
ac & = & k_{1}k_{2}n^{2}+k_{1}dn+k_{2}bn+bd\\
ac & = & bd+n(k_{1}k_{2}n+k_{1}d+k_{2}b)\\
ac-bd & = & n(k_{1}k_{2}n+k_{1}d+k_{2}b)
\end{eqnarray*}
So then $n|(ac-bd)$ and $ac\equiv bd\bmod n$.
\end{proof}
\end{prop}
\begin{xca}
Determine $51\mod 13$, $342\mod 85$, and $(82\cdot73)\mod 7$.
\end{xca}

\begin{xca}
Let $n$ be a positive integer greater than $1$ and let $a$ and
$b$ be positive integers. Prove that $a\mod n\equiv b\mod n$ if
and only if $a\equiv b\mod n$. 
\end{xca}

\section{Mathematical Induction}
\begin{example}
Prove that ${\displaystyle \sum_{k=1}^{n}(2k+1)=n^{2}}$. That is,
prove that $1+3+5+\cdots+(2n-1)=n^{2}$.

\begin{proof}
(step 1: base case) Suppose that $n=1$. Then $1=1^{2}=1$. So then
the statement is true for $n=1$.

(step 2: inductive step) Now \emph{assume that the statement is true
for some $n\geq1$}. We want to show that it is then true for $n+1$.
If it is true for some $n\geq1$ then\emph{ ${\displaystyle \sum_{k=1}^{n+1}}(2k+1)={\displaystyle \sum_{k=1}^{n}(2k+1)+2(n+1)-1}$
}(note how we broke it down into a part we knew and the next step
for $n+1$). 

\emph{
\begin{eqnarray*}
{\displaystyle \sum_{k=1}^{n+1}}(2k+1) & = & {\displaystyle \sum_{k=1}^{n}(2k+1)+2(n+1)-1}=\underset{=n^{2}\mbox{ by the assumption}}{\underbrace{\sum_{k=1}^{n}(2k+1)}}+\underset{=2n+1}{\underbrace{2(n+1)-1}}\\
 & = & n^{2}+2n+1\\
 & = & (n+1)^{2}
\end{eqnarray*}
}so the statement is true by induction.
\end{proof}
\end{example}
Once we have shown it is true for $n=1$ in step $1$ then we can
use step $2$ to show that is true for $n=2$. Once we know it is
true for $n=2$ we know it is true for $n=3$ by step $3$ again.
Continuing like this we can show it is true for any $n$. This process
is called mathematical induction or just proof by induction.
\begin{thm}
\textbf{\textup{The Principle of Mathematical Induction. }}\textup{\emph{Let
$S(n)$ be a statement involving a variable $n$, and the following
hold:}}\emph{}

\begin{enumerate}
\item $S(n_{0})$ is true.
\item If $S(n)$ is true for some positive integer $n\geq n_{0}$ then $S(n+1)$
is true.
\end{enumerate}
Then $S(n)$ is true for all positive integers $n$.
\end{thm}
\begin{proof}
Suppose $S(b)$ is true, and $S(n+1)$ is true whenever $S(n)$ is
true, but there is some $n$ for which $S(n)$ is not true. By the
well ordering principle there is a least number , say $k$, such that
$S(k)$ is not true. $k\neq b$ by assumption, but since $k$ is the
smallest number for which $S$ is not true $S(k-1)$ is true. This
is a contradiction since then $S(k)$ is true by assumption. 
\end{proof}
\begin{xca}
Prove that ${\displaystyle \sum_{k=1}^{n}2k=n^{2}+n}$.
\end{xca}

\begin{xca}
Prove that $23^{n}-1$ is divisible by $11$for all positive integers
$n$.
\end{xca}

\begin{xca}
$\bigstar$ The Fibonacci numbers are :$1,1,2,3,5,8,13,21,34,\dots$.
In general, the Fibonacci numbers are defined by $f_{1}=1$, $f_{2}=1$,
and $f_{n+2}=f_{n+1}+f_{n}$ for $n=1,2,3,\dots$. Prove that the
$n^{\mbox{th}}$ Fibonacci number $f_{n}<2^{n}$.
\end{xca}

\section{Additional problems}
\begin{xca}
$\bigstar$ 

Let $a\neq0$, $b\neq0$, and $d=\gcd\left(a,b\right)$. Prove that
if $d'$ is any common divisor of $a$ and $b$ then $d'|d$.
\end{xca}
\begin{thm}
\textbf{\emph{The Fundamental Theorem of Arithmetic. }}Every integer
greater than 1 is either a prime or the unique product of primes. 
\end{thm}

\begin{xca}
$\bigstar$ If $c^{2}=ab$ and $\gcd\left(a,b\right)=1$, prove that
$a$ and $b$ are perfect squares (Hint: By the theorem above, $a$
and $b$ can be written as the product of primes. They are perfect
squares iff. each prime in this prime factorization has an even power). 

\vspace{0cm}

The set of all congruence classes modulo $n$ is denoted $\mathbb{Z}_{n}$
(read ``$Z$ mod $n$''). 

$\left|\mathbb{Z}_{n}\right|=n$. 
\end{xca}

\begin{xca}
$\bigstar$ Suppose that $a\neq0$, $b\neq0$, and $\mbox{lcm}(a,b)=m$.
Prove that if $n$ is a common multiple of $a$ and $b$ then $m|n$. 
\end{xca}

\begin{xca}
[$\smiley$\textbf{.}]Two sets are said to be \emph{equal }if there
exists a bijective function between them; we write $\left|A\right|=\left|B\right|$.
$B$ is said to have \emph{strictly greater cardinality} than $B$
if there is a one-to-one function from $A$ into $B$, but there does
not exist a bijective function from $A$ into $B$; we write$\left|A\right|<\left|B\right|$.
Prove that $\left|\mathbb{N}\right|<\left|\mathbb{R}\right|$.
\end{xca}

\end{document}
